Integrated circuits are wasteful of energy. Digital circuits charge
transistor gates to change states, and when discharged, the charges are
dumped to ground. In analog circuits the transconductance requires a DC
current, a continuous flow of charges from positive supply to ground.

Integrated circuits are incredibly useful though. Life without would be
different.

A continuous effort from engineers like me have reduced the power
consumption of both digital and analog circuits by order of magnitudes
since the invention of the transistor 75 years ago.

One of the first commercial ADCs, the
\href{https://www.analog.com/media/en/training-seminars/design-handbooks/Data-Conversion-Handbook/Chapter1.pdf}{DATRAC}
on page 24, was a 11-bit 50 kSps that consumed 500 W. That's Walden
figure of merit of 4 \(\mu\)J/conv.step. Today's state-of-the-art ADCs
in the same sampling range have a Walden figure of merit of 0.6
fJ/conv.step {[}1{]}.

4 \(\mu\) / 0.6 f = 6.7e9, a difference in power consumption of almost
10 billion times !!!

Improvements to power consumption have become harder and harder, but I
believe there is still far to go before we cannot reduce power
consumption any more.

Towards a Green and Self-Powered Internet of Things Using Piezoelectric
Energy Harvesting {[}2{]} {[}1{]} has a nice overview of power
consumption of technologies, seen in the next figures below.

In the context of energy harvesting, there is energy in electromagnetic
fields, temperature, and mechanical stress, and there are ways to
translate between them the energy forms.


{
\centering
\aicfig{media/shirv5-2928523-large.png}
}

\emph{Figure 1: The couplings between the electrical, mechanical and
thermal energy domains, including the piezoelectric, pyroelectric and
thermoelectric effects. From Shirvanimoghaddam et al., \emph{IEEE
Access} 7 (2019) {[}2{]}, licensed CC BY 4.0}

Below we can see a figure of the potential energy that can be harvested
per volume, and the type power consumption of technologies {[}1{]}.

As devices approach average power consumption of \(\mu W\) it becomes
possible to harvest the energy from the environment, and do away with
the battery.


{
\centering
\aicfig{media/shirv6-2928523-large.png}
}

\emph{Figure 2: Battery run-time and power consumption of devices from a
32 kHz quartz oscillator to a laptop, plotted against the energy
available from mechanical, thermal and radiant sources in uW/cm2. From
Shirvanimoghaddam et al., \emph{IEEE Access} 7 (2019) {[}2{]}, licensed
CC BY 4.0}

For wireless standards, there are some that can be run on energy
harvesting. Below is an overview from {[}1{]}. Many of us will have a
NFC card in our pocket for payment, or entry to buildings. NFC card has
a integrated circuit that is powered from the electromagnetic field from
the NFC reader.

Other standards, like Bluetooth, WiFi, LTE are harder to run battery
less, because the energy requirement above 1 mW.

Technologies like Bluetooth LE, however, can approach \textless{} 10
\(\mu\)W for some applications, although the burst power may still be 10
mW to 100 mW. As such, although the average power is low, the energy
harvesting cannot support peak loads and a charge storage device is
required (battery, super-capacitor, large capacitor).


{
\centering
\aicfig{media/shirv11-2928523-large.png}
}

\emph{Figure 3: Power versus coverage for wireless standards, showing
that NFC, RFID, Z-Wave and Zigbee can be battery-less while WiFi,
Bluetooth and the cellular standards need a battery. From
Shirvanimoghaddam et al., \emph{IEEE Access} 7 (2019) {[}2{]}, licensed
CC BY 4.0}

I'd like to give you an introduction to the possible ways of harvesting
energy. I know of five methods: - thermoelectric - photovoltaic -
piezoelectric - electromagnetic - triboelectric

\section{\texorpdfstring{\href{https://en.wikipedia.org/wiki/Thermoelectric_effect}{Thermoelectric}}{Thermoelectric}}\label{thermoelectric}

Apply heat to one end of a metal wire, what happens to the free
electrons? As we heat the material we must increase the energy of the
free electrons at the hot end of the wire. The atoms wiggle more, and
when the free electrons scatter off the atomic structure there should be
an exchange of energy. Think of the electrons at the hot side as high
energy electrons, while on the cold side there are low energy electrons,
I think.

There will be diffusion current of electrons in both directions in the
material, however, if the mobility of electrons in the material is
dependent on the energy, then we would get a difference in current of
low energy electrons and high energy electrons. A difference in current
would lead to a charge difference at the hot end and cold end, which
would give a difference in voltage.

Take a copper wire, bend it in half, heat the end with the loop, and
measure the voltage at the cold end. Would we measure a voltage
difference?

\textbf{NO}, there would not be a voltage difference between the two
ends of the wire. The voltage on the loop side would be different, but
on the cold side, where we have the ends, there would be no voltage
difference.

Gauss law tell us that inside a conductor there cannot be a static field
without a current. As such, if there was a voltage difference between
the cold ends, it would quickly be dissipated, and no DC current would
flow.

The voltage difference in the material between the hot and cold end will
create currents, but we can't use them if we only have one type of
material.

Imagine we have Iron and copper wires, as shown below, and we heat one
end. In that case, we can draw current between the cold ends.


{
\centering
\aicfig{media/Thermoelectric_effect.png}
}

\emph{Figure 4: The thermoelectric effect - iron and copper wires joined
at a heated end drive a current through a meter connected at the cold
end. Image: Cmglee, CC BY-SA 4.0, via Wikimedia Commons}

The voltage difference at the hot and cold end is described by the

\href{https://en.wikipedia.org/wiki/Seebeck_coefficient}{Seebeck
coefficient}

Imagine two parallel wires with different Seebeck coefficients, one of
copper (\(6.5\text{ } \mu V/K\)) and one of iron
(\(19\text{ } \mu V/K\)). We connect them at the hot end. The voltage
difference between hot and cold would be higher in the iron, than in the
copper. At the cold end, we would now measure a difference in voltage
between the wires!

In silicon, the Seebeck coefficient can be modified through doping. A
model of Seebeck coefficient is shown below. The value of the Seebeck
coefficient depends on the location of the Fermi level in relation to
the Conduction band or the valence band.


{
\centering
\aicfig{media/Absolute_Seebeck_coefficients_of_various_metals_up_to_high_temperatures.png}
}

\emph{Figure 5: Absolute Seebeck coefficient versus temperature for a
range of metals, spanning roughly +20 uV/K for tungsten to -60 uV/K for
palladium. Image: Nanite, CC0, via Wikimedia Commons}


{
\centering
\aicfig{media/Mott_Seebeck_silicon.png}
}

\emph{Figure 6: Seebeck coefficient and conductivity of silicon as a
function of Fermi level, with the sign flipping between the valence and
conduction band edges. Image: Nanite, CC0, via Wikimedia Commons}

In the picture below we have a silicon (the cyan and yellow colors).

Assume we dope with acceptors (yellow, p-type), that shifts the Fermi
level closer to the Valence band (\(E_V\)), and the dominant current
transport will be by holes, maybe we get 1 mV/K from the picture above.

For the material doped with donors (cyan, n-type) the Fermi level is
shifted towards the Conduction band (\(E_C\)), and the dominant charge
transport is by electrons, maybe we get -1 mV/K from the picture above.

Assume we have a temperature difference of 50 degrees, then maybe we
could get a voltage difference at the cold end of 100 mV. That's a low
voltage, but is possible to use.


{
\centering
\aicfig{media/Thermoelectric_Generator_Diagram.png}
}

\emph{Figure 7: Thermoelectric generator, where n-type and p-type legs
between a heat source and a cool side drive a current through the load
resistor. Image: Ken Brazier, CC BY-SA 4.0, via Wikimedia Commons}

The process can be run in reverse. In the picture below we force a
current through the material, we heat one end, and cool the other. Maybe
you've heard of
\href{https://en.wikipedia.org/wiki/Thermoelectric_cooling}{Peltier
elements}.


{
\centering
\aicfig{media/Thermoelectric_Cooler_Diagram.png}
}

\emph{Figure 8: The same structure run in reverse as a Peltier cooler,
where a forced current cools the top surface and dissipates heat at the
bottom. Image: Ken Brazier, CC BY-SA 4.0, via Wikimedia Commons}

\subsection{\texorpdfstring{\href{https://en.wikipedia.org/wiki/Radioisotope_thermoelectric_generator}{Radioisotope
Thermoelectric
generator}}{Radioisotope Thermoelectric generator}}\label{radioisotope-thermoelectric-generator}

Maybe you've heard of a nuclear battery. Sounds fancy, right? Must be
complicated, right?

Not really, take some radioactive material, which generates heat, stick
a thermoelectric generator to the hot side, make sure you can cool the
cold side, and we have a nuclear battery.

Nuclear batteries are ``simple'', and ultra reliable. There's not really
a chemical reaction. The nucleus of the radioactive material degrades,
but not fast. In the thermoelectric generator, there are no moving
parts.

In a normal battery there is a chemical reaction that happens when we
pull a current. Atoms move around. Eventually the chemical battery will
change and degrade.

Nuclear batteries were used in Voyager, and they still work to this day.
The nuclear battery is the round thing underneath Voyager in the picture
below. The radioisotopes provide the heat, space provides the cold, and
voila, \href{https://en.wikipedia.org/wiki/Voyager_1}{470 W} to run the
electronics.


{
\centering
\aicfig{media/1280px-Voyager_spacecraft_model.png}
}

\emph{Figure 9: The Voyager spacecraft, whose radioisotope
thermoelectric generator is the finned cylinder below the dish. Image:
NASA, public domain, via Wikimedia Commons}

\subsection{\texorpdfstring{\href{https://en.wikipedia.org/wiki/Thermoelectric_generator}{Thermoelectric
generators}}{Thermoelectric generators}}\label{thermoelectric-generators}

Assume we wanted to drive a watch from a thermoelectric generator (TEG).
The skin temperature is maybe 33 degrees Celsius, while the ambient
temperature is maybe 23 degrees Celsius on average.

From the model of a thermoelectric generator below we'd get a voltage of
10 mV to 500 mV, too low for most integrated circuits.

In order to drive an integrated circuit we'd need to boost the voltage
to maybe 1.8 V.

The main challenge with thermoelectric generators is to provide a
cold-boot function where the energy harvester starts up at a low
voltage.

In silicon, it is tricky to make anything work below some thermal
voltages (kT/q). We at least need about 3 -- 4 thermal voltages to make
anything function.

The key enabler for an efficient, low temperature differential, energy
harvester is an oscillator that works at low voltage (i.e 75 mV). If we
have a clock, then we can boost with capacitors

In A 3.5-mV Input Single-Inductor Self-Starting Boost Converter With
Loss-Aware MPPT for Efficient Autonomous Body-Heat Energy Harvesting
{[}3{]} they use a combination of both switched capacitor and switched
inductor boost.


{
\centering
\aicfig{media/l11_teg_mdl_tikz.png}
}

\emph{Figure 10: Model of a thermoelectric generator - a 1 to 50 mV/K
source in series with a source resistance below 10 ohm - together with
the self-starting boost converter architecture from Bose et al.}

\section{\texorpdfstring{\href{https://en.wikipedia.org/wiki/Photovoltaic_effect}{Photovoltaic}}{Photovoltaic}}\label{photovoltaic}

In silicon, photons can knock out electron/hole pairs. If we have a PN
junction, then it's possible to separate the electron/holes before they
recombine as shown in figure below.

An electron/hole pair knocked out in the depletion region (1) will
separate due to the built-in field. The hole will go to P and the
electron to N. This increases the voltage VD across the diode.

A similar effect will occur if the electron/hole pair is knocked out in
the P region (2). Although the P region has an abundance of holes, the
electron will not recombine immediately. If the electron diffuses close
to the depletion region, then it will be swept across to the N side, and
further increase VD.

On the N-side the same minority carrier effect would further increase
the voltage (3).


{
\centering
\aicfig{media/l11_pv_pn_tikz.png}
}

\emph{Figure 11: Photon absorption in a PN junction, where electron/hole
pairs generated in the depletion region (1), the P region (2) and the N
region (3) all increase the diode voltage VD}

A circuit model of a Photodiode can be seen in figure below, where it is
assumed that a single photodiode is used. It is possible to stack
photodiodes to get a higher output voltage.


{
\centering
\aicfig{media/l11_pv_mdl_tikz.png}
}

\emph{Figure 12: Circuit model of a photodiode - a photo current source
in parallel with a diode, delivering a load current IL}

As the load current is increased, the voltage VD will drop. As the photo
current is increased, the voltage VD will increase. As such, there is an
optimum current load where there is a balance between the photocurrent,
the voltage VD and the load current.

\[ I_D = I_S\left(e^\frac{V_D}{V_T} - 1\right)\]

\[ I_D = I_{Photo} - I_{Load}\]

\[ V_D = V_T ln{\left(\frac{I_{Photo} - I_{Load}}{I_S} + 1 \right)} \]

\[ P_{Load} = V_D I_{Load}\]

Below is a model of the power in the load as a function of diode voltage

\begin{Shaded}
\begin{Highlighting}[]
\CommentTok{\#!/usr/bin/env python3}
\ImportTok{import}\NormalTok{ numpy }\ImportTok{as}\NormalTok{ np}
\ImportTok{import}\NormalTok{ matplotlib.pyplot }\ImportTok{as}\NormalTok{ plt}

\NormalTok{m }\OperatorTok{=} \FloatTok{1e{-}3}  
\NormalTok{i\_load }\OperatorTok{=}\NormalTok{ np.linspace(}\FloatTok{1e{-}5}\NormalTok{,}\FloatTok{1e{-}3}\NormalTok{,}\DecValTok{200}\NormalTok{)}

\NormalTok{i\_s }\OperatorTok{=} \FloatTok{1e{-}12}  \CommentTok{\# saturation current }
\NormalTok{i\_ph }\OperatorTok{=} \FloatTok{1e{-}3}  \CommentTok{\# Photocurrent}

\NormalTok{V\_T }\OperatorTok{=} \FloatTok{1.38e{-}23}\OperatorTok{*}\DecValTok{300}\OperatorTok{/}\FloatTok{1.6e{-}19}  \CommentTok{\#Thermal voltage}

\NormalTok{V\_D }\OperatorTok{=}\NormalTok{ V\_T}\OperatorTok{*}\NormalTok{np.log((i\_ph }\OperatorTok{{-}}\NormalTok{ i\_load)}\OperatorTok{/}\NormalTok{(i\_s) }\OperatorTok{+} \DecValTok{1}\NormalTok{)}

\NormalTok{P\_load }\OperatorTok{=}\NormalTok{ V\_D}\OperatorTok{*}\NormalTok{i\_load}

\NormalTok{plt.subplot(}\DecValTok{2}\NormalTok{,}\DecValTok{1}\NormalTok{,}\DecValTok{1}\NormalTok{)}
\NormalTok{plt.plot(i\_load}\OperatorTok{/}\NormalTok{m,V\_D)}
\NormalTok{plt.ylabel(}\StringTok{"Diode voltage [mA]"}\NormalTok{)}
\NormalTok{plt.grid()}
\NormalTok{plt.subplot(}\DecValTok{2}\NormalTok{,}\DecValTok{1}\NormalTok{,}\DecValTok{2}\NormalTok{)}
\NormalTok{plt.plot(i\_load}\OperatorTok{/}\NormalTok{m,P\_load}\OperatorTok{/}\NormalTok{m)}
\NormalTok{plt.xlabel(}\StringTok{"Current load [mA]"}\NormalTok{)}
\NormalTok{plt.ylabel(}\StringTok{"Power Load [mW]"}\NormalTok{)}
\NormalTok{plt.grid()}
\NormalTok{plt.savefig(}\StringTok{"pv.pdf"}\NormalTok{)}
\NormalTok{plt.show()}
\end{Highlighting}
\end{Shaded}

From the plot below we can see that to optimize the power we could
extract from the photovoltaic cell we'd want to have a current of 0.9 mA
in the model above.

The
\href{https://wulffern.github.io/aic2026/assets/examples/pv.html}{interactive
version of this example} plots the same equation both ways round, marks
the maximum power point and reports the fill factor. Sweeping the
photocurrent over six decades shows why harvesting in a dim room is a
current problem rather than a voltage one.


{
\centering
\aicfig{media/pv_tikz.png}
}

\emph{Figure 13: Diode voltage and delivered power versus load current
for the photodiode model, with the maximum power point near 0.9 mA}

Most photovoltaic energy harvesting circuits will include a maximum
power point tracker as the optimum changes with light conditions.

In A Reconfigurable Capacitive Power Converter With Capacitance
Redistribution for Indoor Light-Powered Batteryless Internet-of-Things
Devices {[}4{]} they include a maximum power point tracker and a
reconfigurable charge pump to optimize efficiency.

\section{\texorpdfstring{\href{https://en.wikipedia.org/wiki/Piezoelectricity}{Piezoelectric}}{Piezoelectric}}\label{piezoelectric}

I'm not sure I understand the piezoelectric effect, but I think it goes
something like this.

Consider a crystal made of a combination of elements, for example
\href{http://lampx.tugraz.at/~hadley/ss1/crystalstructure/structures/semiconductors/GaN.html}{Gallium
Nitride}. In GaN it's possible to get a polarization of the unit cell,
with a more negative charge on one side, and a positive charge on the
other side. The polarization comes from an asymmetry in the electron and
nucleus distribution within the material.

In a polycrystaline substance the polarization domains will usually be
random, and no electric field will observable. The polarization domains
can be aligned by heating the material and applying a electric field.
Now all the small electric fields point in the same direction.

From Gausses law we know that the electric field through a surface is
determined by the volume integral of the charges inside.

\[ \oint_{\partial \Omega} \mathbf{E} \cdot d\mathbf{S} = \frac{1}{\epsilon_0} \iiint_{V} \rho
\cdot dV\]

Although there is a net zero charge inside the material, there is an
uneven distribution of charges, as such, some of the field lines will
cross through the surface.

Assume we have a polycrystaline GaN material with polarized domains. If
we measure the voltage across the material we will read 0 V. Even though
the domains are polarized, and we should observe an external electric
field, the free charges in the material will redistribute if there is a
field inside, such that there is no current flowing, and thus no
external field.

If we apply stress, however, all the domains inside the material will
shift. Now the free charges do not exactly cancel the electric field in
the material, the free charges are in the wrong place. If we have a
material with low conductivity, then it will take time for the free
charges to redistribute. As such, for a while, we can measure an voltage
across the material.

Assuming the above explanation is true, then there should not be
piezoelectric materials with high conductivity, and indeed, most
piezoelectric materials have resistance of
\href{https://www.f3lix-tutorial.com/piezo-materials}{\(10^{12}\) to
\(10^{14}\) Ohm}.

Vibrations on a piezoelectric material will result in a AC voltage
across the surface, which we can harvest.

A model of a piezoelectric transducer can be seen below.

The voltage on the transducer can be on the order of a few volts, but
the current is usually low (nA -- µA). The key challenge is to rectify
the AC signal into a DC signal. It is common to use tricks to reduce the
energy waste due to the rectifier.

An example of piezoelectric energy harvester can be found in A Fully
Integrated Split-Electrode SSHC Rectifier for Piezoelectric Energy
Harvesting {[}5{]}


{
\centering
\aicfig{media/lx_piezo_mdl_tikz.png}
}

\emph{Figure 14: Model of a piezoelectric transducer - an AC current
source in parallel with the transducer capacitance}

\section{Electromagnetic}\label{electromagnetic}\index{electromagnetic@Electromagnetic}

\subsection{``Near field'' harvesting}\label{near-field-harvesting}

Near Field Communication (NFC) operates at close physical distances


{
\centering
\aicfig{media/FarNearFields-USP-4998112-1.png}
}

\emph{Figure 15: The reactive near field, radiative (Fresnel) near field
and far field regions around an antenna. Image: Goran M Djuknic, public
domain (US patent), via Wikimedia Commons}

Reactive near field or inductive near field

\[ \text{Inductive} < \frac{\lambda}{2 \pi}\]

Within the inductive near field the antenna's can ``feel'' each other.
The NFC reader inside the card reader can ``feel'' the antenna of the
NFC tag. When the tag get's close it will load down the NFC reader by
presenting a load impedance. As the circuit inside the tag is powered,
it can change the impedance of it's antenna, which is sensed by the
reader, and thus the reader can get data from the tag. The tag could
lock in on the 13.56 MHz frequency and decode both amplitude and phase
modulation from the reader.

Since the NFC or Qi system operates at close distances, then the
coupling factor between antenna's, or really, inductors, can be decent,
and it's possible to achieve efficiencies of maybe 70 \%.

At Bluetooth frequencies, as can be seen below, it does not really make
sense to couple inductors, as they need to be within 2 cm to be in the
inductive near field. The inductive near field is a significant problem
for the coupling between inductors on chip, but I don't think I would
use it to transfer power.

\begin{longtable}[]{@{}lcc@{}}
\toprule\noalign{}
Standard & Frequency {[}MHz{]} & Inductive {[}m{]} \\
\midrule\noalign{}
\endhead
\bottomrule\noalign{}
\endlastfoot
AirFuel Resonant & 6.78 & 7.03 \\
NFC & 13.56 & 3.52 \\
Qi & 0.205 & 232 \\
Bluetooth & 2400 & 0.02 \\
\end{longtable}

\subsection{Ambient RF Harvesting}\label{ambient-rf-harvesting}\index{ambient rf harvesting@Ambient RF Harvesting}

Extremely inefficient idea, but may find special use-cases at
short-distance.

Will get better with beam-forming and directive antennas

There are companies that think RF harvesting is a good idea.

\href{https://airfuel.org/airfuel-rf/}{AirFuel RF}

I think that ambient RF harvesting should tingle your science spidy
senses.

Let's consider the power transmitted in wireless standards. Your
cellphone may transmit 30 dBm, your WiFi router maybe 20 dBm, and your
Bluetooth LE device 10 dBm.

In case those numbers don't mean anything to you, below is a conversion
to watts.

\begin{longtable}[]{@{}cc@{}}
\toprule\noalign{}
dBm & W \\
\midrule\noalign{}
\endhead
\bottomrule\noalign{}
\endlastfoot
30 & 1 \\
0 & 1 m \\
-30 & 1 u \\
-60 & 1 n \\
-90 & 1 p \\
\end{longtable}

Now ask your self the question ``What's the power at a certain
distance?''. It's easier to flip the question, and use Friis to
calculate the distance.

Assume \[P_{TX}\] = 1 W (30 dBm) and \[P_{RX}\] = 10 uW (-20 dBm)

then

\[ D = 10^\frac{P_{TX} - P_{RX} + 20 log_{10}\left(\frac{c}{4 \pi f}\right)}{20} \]

In the table below we can see the distance is not that far!

\begin{longtable}[]{@{}lcr@{}}
\toprule\noalign{}
Freq & \textbf{\(20 log_{10}\left(c/4 \pi f\right)\)} {[}dB{]} & D
{[}m{]} \\
\midrule\noalign{}
\endhead
\bottomrule\noalign{}
\endlastfoot
915M & -31.7 & 8.2 \\
2.45G & -40.2 & 3.1 \\
5.80G & -47.7 & 1.3 \\
\end{longtable}

I believe ambient RF is a stupid idea.

Assuming an antenna that transmits equally in all direction, then the
loss on the first meter is 40 dB at 2.4 GHz. If I transmitted 1 W, there
would only be 100 µW available at 1 meter. That's an efficiency of 0.01
\%.

Just fundamentally stupid. \textbf{Stupid, I tell you!!!}

\emph{Stupidity in engineering really annoys me, especially when people
don't understand how stupid ideas are.}

\section{Triboelectric generator}\label{triboelectric-generator}\index{triboelectric generator@Triboelectric generator}

Although static electricity is an old phenomenon, it is only recently
that triboelectric nanogenerators have been used to harvest energy.

An overview can be seen in Current progress on power management systems
for triboelectric nanogenerators {[}6{]}.

A model of a triboelectric generator can be seen in below. Although the
current is low (nA) the voltage can be high, tens to hundreds of volts.

The key circuit challenge is the rectifier, and the high voltage output
of the triboelectric generator.

Take a look in A Fully Energy-Autonomous Temperature-to-Time Converter
Powered by a Triboelectric Energy Harvester for Biomedical Applications
{[}7{]} for more details.


{
\centering
\aicfig{media/lx_trib_mdl_tikz.png}
}

\emph{Figure 16: Model of a triboelectric generator - an AC source in
series with the transducer capacitance}

Tan et al. {[}7{]} built exactly this, and it is worth walking through
because it is a complete system rather than a circuit.

Their transducer is a square centimetre of conductive nickel against
nickel-PFA, and the energy source is human motion below 1 Hz. So the
input is not a waveform in any useful sense: it is a few sparse pulses a
second, each of them high voltage and almost no current, and between
them nothing at all. The model in the figure above is the right picture
for it, with one addition - a real transducer also has a shunt
resistance across it, which is one more path for the little charge there
is to leak away before it can be used.

That shapes the whole design. The rectifier cannot be a diode bridge,
because two diode drops of the little charge available leaves nothing,
and it cannot leak, because the next pulse may be a second away. Their
answer is a low leakage rectifier feeding a power management unit whose
static consumption is measured in nanowatts, storing charge on a
capacitor until there is enough to do something with. On their
measurements a conventional rectifier never reaches the 600 mV the
circuit needs to start; theirs does.

What it then does is neat. Rather than digitise a temperature, it starts
a PTAT bandgap, charges a capacitor and waits for a comparator to trip,
so the output is a pulse whose \emph{width} is the temperature. No
clock, no converter, no reference to speak of - the temperature falls
out of a single one-shot discharge. If all you have is a millijoule now
and then, answering in the time domain is a great deal cheaper than
answering in the voltage domain.

One detail in their block diagram is worth more than the circuit though.
There is a \texttt{VDD\_ext} on it. The system is not fully harvested.

That is not a criticism, and it is a good illustration of the division
of labour. They set out to show that a triboelectric harvester can run a
temperature sensor, and they showed it; how the reading gets from the
chip to wherever anyone would read it is left alone, and getting a radio
to run off sub-1 Hz human motion is a different paper. It is academia's
job to prove that something could be possible, and industry's job to
make something that could be possible actually work. Read papers with
that in mind and look for the block that is still plugged into the wall.

\section{Comparison}\label{comparison}

Imagine you're a engineer in a company that makes integrated circuits.
Your CEO comes to you and says ``You need to make a power management IC
that harvest energy and works with everything''.

Hopefully, your response would now be ``That's a stupid idea, any energy
harvester circuit must be made specifically for the energy source''.

Thermoelectric and photovoltaic provide low DC voltage. Piezoelectric
and Triboelectric provide an AC voltage. Ambient RF is stupid.

For a ``energy harvesting circuit'' you must also know the application
(wrist watch, or wall switch) to know what energy source is available.

Below is a table that show's a comparison in the power that can be
extracted.

The power levels below are too low for the peak power consumption of
integrated circuits, so most applications must include a charge storage
device, either a battery, or a capacitor.

\begin{longtable}[]{@{}
  >{\raggedright\arraybackslash}p{(\linewidth - 6\tabcolsep) * \real{0.1679}}
  >{\raggedright\arraybackslash}p{(\linewidth - 6\tabcolsep) * \real{0.3893}}
  >{\raggedright\arraybackslash}p{(\linewidth - 6\tabcolsep) * \real{0.1069}}
  >{\raggedright\arraybackslash}p{(\linewidth - 6\tabcolsep) * \real{0.3359}}@{}}
\toprule\noalign{}
\begin{minipage}[b]{\linewidth}\raggedright
Energy source
\end{minipage} & \begin{minipage}[b]{\linewidth}\raggedright
Power density
\end{minipage} & \begin{minipage}[b]{\linewidth}\raggedright
Frequency
\end{minipage} & \begin{minipage}[b]{\linewidth}\raggedright
Characteristics
\end{minipage} \\
\midrule\noalign{}
\endhead
\bottomrule\noalign{}
\endlastfoot
Solar / PV & 10 uW/cm\(^2\) indoor, 15 mW/cm\(^2\) outdoor & DC &
Requires exposure to light \\
RF & 0.1 uW/cm\(^2\) GSM, 0.01 uW/cm\(^2\) WiFi & 380 MHz--5 GHz & Poor
indoors and out of line of sight \\
Thermal, body heat & 40 uW/cm\(^2\) & DC & Requires a high temperature
difference \\
Piezoelectric & 4 uW/cm\(^2\) & \textgreater{} 30 Hz & Not limited to
indoors or outdoors \\
Triboelectric (TENG) & 1 uW/cm\(^2\) & 1 Hz & Not limited to indoors or
outdoors \\
\end{longtable}

Numbers from Tan et al. {[}7{]}.

Read the first column against the last. Solar wins outdoors by three
orders of magnitude and loses most of that the moment you go inside, RF
is worse than either everywhere, and the two mechanical sources are the
only ones that do not care where they are. That is the whole argument
for bothering with piezoelectric and triboelectric harvesting despite
their being at the bottom of the power column: a microwatt you can rely
on beats a milliwatt that depends on someone opening the curtains.

The frequency column matters as much as the power column and is easier
to overlook. DC sources want a boost converter; the mechanical ones
deliver at a few hertz or less and want a rectifier and a reservoir, and
at 1 Hz the circuit spends almost all of its life waiting. That is why
the leakage of the storage path, rather than the efficiency of the
conversion, is what decides whether a triboelectric system works.

\section{Summary}\label{summary}

The one-page version of this chapter:

\begin{itemize}
\tightlist
\item
  Harvesters deliver power, batteries deliver energy: a harvester design
  starts from the average load current, not the peak
\item
  Thermoelectric generators give millivolts per kelvin of gradient -
  real designs start below 100 mV and need a boost converter that can
  cold-start from that little
\item
  A photovoltaic cell is the photodiode from this chapter run in the
  fourth quadrant: half a volt a cell, current proportional to light
\item
  Piezo and electromagnetic harvesters turn vibration into AC that must
  be rectified; triboelectric is the same story with contact charging
\item
  Ambient RF sounds free until Friis has spoken: microwatts at best, and
  only near the transmitter
\item
  Every source needs power management sized to its impedance - maximum
  power transfer is the whole game
\end{itemize}

\section{Would you like to know
more?}\label{would-you-like-to-know-more}

{[}1{]} Towards a Green and Self-Powered Internet of Things Using
Piezoelectric Energy Harvesting {[}2{]}

A 3.5-mV Input Single-Inductor Self-Starting Boost Converter With
Loss-Aware MPPT for Efficient Autonomous Body-Heat Energy Harvesting
{[}3{]}

A Reconfigurable Capacitive Power Converter With Capacitance
Redistribution for Indoor Light-Powered Batteryless Internet- of-Things
Devices {[}4{]}

A Fully Integrated Split-Electrode SSHC Rectifier for Piezoelectric
Energy Harvesting {[}5{]}

Current progress on power management systems for triboelectric
nanogenerators {[}6{]}

A Fully Energy-Autonomous Temperature-to-Time Converter Powered by a
Triboelectric Energy Harvester for Biomedical Applications {[}7{]}

\phantomsection\label{refs}
\begin{CSLReferences}{0}{0}
\bibitem[\citeproctext]{ref-hsieh18}
\CSLLeftMargin{{[}1{]} }%
\CSLRightInline{S.-E. Hsieh and C.-C. Hsieh, {``A 0.4V 13b 270kS/s
SAR-ISDM ADC with an opamp-less time-domain integrator,''} in \emph{2018
IEEE international solid-state circuits conference - (ISSCC)}, 2018, pp.
240--242, doi:
\href{https://doi.org/10.1109/ISSCC.2018.8310273}{10.1109/ISSCC.2018.8310273}.
}

\bibitem[\citeproctext]{ref-shirvanimoghaddam19}
\CSLLeftMargin{{[}2{]} }%
\CSLRightInline{M. Shirvanimoghaddam \emph{et al.}, {``Towards a green
and self-powered internet of things using piezoelectric energy
harvesting,''} \emph{IEEE Access}, vol. 7, pp. 94533--94556, 2019, doi:
\href{https://doi.org/10.1109/ACCESS.2019.2928523}{10.1109/ACCESS.2019.2928523}.
}

\bibitem[\citeproctext]{ref-bose21}
\CSLLeftMargin{{[}3{]} }%
\CSLRightInline{S. Bose, T. Anand, and M. L. Johnston, {``A 3.5-mV input
single-inductor self-starting boost converter with loss-aware MPPT for
efficient autonomous body-heat energy harvesting,''} \emph{IEEE Journal
of Solid-State Circuits}, vol. 56, no. 6, pp. 1837--1848, 2021, doi:
\href{https://doi.org/10.1109/JSSC.2020.3042962}{10.1109/JSSC.2020.3042962}.
}

\bibitem[\citeproctext]{ref-cheng21}
\CSLLeftMargin{{[}4{]} }%
\CSLRightInline{H.-C. Cheng, P.-H. Chen, Y.-T. Su, and P.-H. Chen, {``A
reconfigurable capacitive power converter with capacitance
redistribution for indoor light-powered batteryless internet-of-things
devices,''} \emph{IEEE Journal of Solid-State Circuits}, vol. 56, no.
10, pp. 2934--2942, 2021, doi:
\href{https://doi.org/10.1109/JSSC.2021.3075217}{10.1109/JSSC.2021.3075217}.
}

\bibitem[\citeproctext]{ref-du19}
\CSLLeftMargin{{[}5{]} }%
\CSLRightInline{S. Du, Y. Jia, C. Zhao, G. A. J. Amaratunga, and A. A.
Seshia, {``A fully integrated split-electrode SSHC rectifier for
piezoelectric energy harvesting,''} \emph{IEEE Journal of Solid-State
Circuits}, vol. 54, no. 6, pp. 1733--1743, 2019, doi:
\href{https://doi.org/10.1109/JSSC.2019.2893525}{10.1109/JSSC.2019.2893525}.
}

\bibitem[\citeproctext]{ref-hu22}
\CSLLeftMargin{{[}6{]} }%
\CSLRightInline{T. Hu, H. Wang, W. Harmon, D. Bamgboje, and Z.-L. Wang,
{``Current progress on power management systems for triboelectric
nanogenerators,''} \emph{IEEE Transactions on Power Electronics}, vol.
37, no. 8, pp. 9850--9864, 2022, doi:
\href{https://doi.org/10.1109/TPEL.2022.3156871}{10.1109/TPEL.2022.3156871}.
}

\bibitem[\citeproctext]{ref-tan21}
\CSLLeftMargin{{[}7{]} }%
\CSLRightInline{J. S. Y. Tan \emph{et al.}, {``A fully energy-autonomous
temperature-to-time converter powered by a triboelectric energy
harvester for biomedical applications,''} \emph{IEEE Journal of
Solid-State Circuits}, vol. 56, no. 10, pp. 2913--2923, 2021, doi:
\href{https://doi.org/10.1109/JSSC.2021.3080383}{10.1109/JSSC.2021.3080383}.
}

\end{CSLReferences}
